\practicechapter{第 3 章实践：GEMM、packing 与 SIMD 证据}

本章对应原书中的第三册“从矩阵向量乘走向 GEMM、packing 与 SIMD”。不要把它当作概念复述，本章要把 packed weight、scalar baseline、AVX2 backend、output block、shape sweep、反汇编和降级 放到同一个可运行项目里检查。贯穿代码仍然是 \filepath{books/cpu-volume-3/labs/linux\_cpu\_inference}；如果某个实验在当前机器上不能验证，报告必须写清降级原因，不能把源码存在当作实测证据。

\practicesection{一、对应原书问题：概念必须落到对象和命令}

原书讲的是系统边界，本章实践要回答三个问题。第一，哪个代码对象承载这个概念；第二，哪个命令能产生可复查输出；第三，哪个坏输入或缺失字段能证明边界不是空话。这里的核心对象包括 \code{output_block}，核心命令或测试入口是 \cmd{lcqi_bench}。

\begin{lstlisting}[numbers=none,caption={第 3 章实践：GEMM、packing 与 SIMD 证据 的最小命令入口}]
cmake -S books/cpu-volume-3 -B books/cpu-volume-3/build/lcqi-debug -DCMAKE_BUILD_TYPE=Debug
cmake --build books/cpu-volume-3/build/lcqi-debug
ctest --test-dir books/cpu-volume-3/build/lcqi-debug --output-on-failure
# chapter focus: lcqi_bench / output_block
\end{lstlisting}

\practicesection{二、从特例推到规律：先抓一个能手算的切片}

本章先看特例：同一组输入同时跑 \code{scalar}、\code{packed\_scalar} 和 \code{avx2}。读者要先手写预期结果，再运行项目观察输出。动作是：比较 CSV 中 \code{backend}、\code{available}、\code{average\_us}、\code{max\_abs\_diff} 和 \code{checksum}。如果输出里没有 \code{max\_abs\_diff} 或对应字段无法解释，就不要继续优化；先回到 reference、输入合同和 trace 字段。

从这个特例推出的规律是：SIMD 路径只有在结果和 baseline 对齐后才有资格谈速度；\code{available=0} 也是证据，不是失败 这条规律要写进报告，不要只写“运行成功”。运行成功只说明某条路径没崩，解释清楚输入、输出、错误路径和不可验证边界，才说明学会了这个知识点。

\practicesection{三、实践任务：把观察固定成报告字段}

任务一：定位源码。至少阅读 \filepath{books/cpu-volume-3/labs/linux\_cpu\_inference/README.md}、相关 \filepath{include/lcqi} 头文件和一个 \filepath{src} 实现文件。报告要写出函数入口、输入对象、输出对象和错误处理方式。

任务二：运行命令。保存构建命令、测试命令、核心输出和环境字段。若命令依赖 AVX2、perf、GPU、NUMA 或外部库，必须记录当前机器是否支持；不支持时把结论降级为“接口已设计，性能未验证”。

任务三：构造反例。围绕本章知识点设计一个坏输入、缺失字段或不满足前提的场景，说明系统应拒绝、降级、fallback 还是继续执行。只给正常路径不合格。

\practicesection{四、验收和递进大作业}

验收标准：报告能从一个具体输入推到工程规则；命令可以在仓库中复跑；输出字段能对应到源码；失败路径说得清；未验证边界不伪装成结论。

递进大作业：提交一个 shape sweep 表，至少解释一个 SIMD 有收益、一个收益不明显、一个当前环境不可验证的场景。这个作业会被后续章节继续引用，命名、目录和字段不要随意变化。
